\section{总结与延伸}

\subsection{讲者的核心总结}

\begin{figure}[H]
\centering
\includegraphics[width=0.95\textwidth]{slides-images/slide-085.jpg}
\caption{讲座结尾：Nathan Lambert 的联系方式与致谢\protect\footnotemark}
\end{figure}
\footnotetext{Slides 第86页。}

Nathan Lambert 在讲座中传递了几个核心信息：

\begin{enumerate}
    \item \textbf{DPO 是正确的起点}：对于任何想进入对齐研究的人，DPO 简单、有效、易于调试
    \item \textbf{数据比算法更重要}：UltraFeedback 数据集的影响力远超任何单一算法改进
    \item \textbf{PPO 略优但代价高昂}：如果你有资源，PPO 能给你额外 1\% 的提升
    \item \textbf{在线方法是未来}：无论是 PPO 的天然在线性，还是 DPO 的在线变体，数据新鲜度都至关重要
    \item \textbf{评估是基础设施}：没有好的评估工具，我们就在黑暗中摸索
\end{enumerate}

\subsection{全课知识图谱}

\begin{center}
\begin{tikzpicture}[node distance=0.9cm, every node/.style={draw, rounded corners, minimum width=3.5cm, minimum height=0.7cm, font=\small}]
\node (pretrain) {预训练语言模型};
\node (sft) [below=of pretrain] {指令微调（SFT/IFT）};
\node (pref) [below=of sft] {偏好优化（DPO/PPO）};
\node (eval) [below left=1cm and -0.5cm of pref] {评估（RewardBench）};
\node (online) [below right=1cm and -0.5cm of pref] {在线方法};
\node (future) [below=2.5cm of pref] {下一代对齐};
\draw[->, thick] (pretrain) -- (sft);
\draw[->, thick] (sft) -- (pref);
\draw[->, thick] (pref) -- (eval);
\draw[->, thick] (pref) -- (online);
\draw[->, thick] (eval) -- (future);
\draw[->, thick] (online) -- (future);
\node[draw=none, right=1.5cm of pretrain, text width=4.5cm, font=\scriptsize] {Transformer, GPT, Scaling Laws};
\node[draw=none, right=1.5cm of sft, text width=4.5cm, font=\scriptsize] {Alpaca, Vicuna, ShareGPT, UltraFeedback};
\node[draw=none, right=1.5cm of pref, text width=4.5cm, font=\scriptsize] {Zephyr, T\"ulu 2, Llama 3};
\end{tikzpicture}
\end{center}

\subsection{关键 Takeaways}

\begin{importantbox}{六条核心原则}
\begin{enumerate}
    \item \textbf{RLHF 是必要但不充分的}：预训练提供基础，后训练使模型"可用"
    \item \textbf{从 DPO 开始}：简单、有效、可扩展，是任何对齐项目的正确起点
    \item \textbf{数据决定上限}：算法只能在数据质量允许的范围内优化
    \item \textbf{评估驱动研发}：没有好的评估工具，就无法判断改进是否有效
    \item \textbf{在线 > 离线}：新鲜的、来自当前策略的数据比历史数据更有价值
    \item \textbf{对齐是多阶段过程}：工业界在不同阶段组合使用 SFT、Rejection Sampling、DPO 和 PPO
\end{enumerate}
\end{importantbox}

\subsection{拓展阅读}

\begin{itemize}
    \item Rafailov et al., \textit{Direct Preference Optimization: Your Language Model is Secretly a Reward Model}, 2023: \url{https://arxiv.org/abs/2305.18290}
    \item Lambert et al., \textit{RewardBench: Evaluating Reward Models for Language Modeling}, 2024: \url{https://arxiv.org/abs/2403.13787}
    \item Ivison et al., \textit{Unpacking DPO and PPO: Disentangling Best Practices for Learning from Preference Feedback}, 2024
    \item Tunstall et al., \textit{Zephyr: Direct Distillation of LM Alignment}, 2023: \url{https://arxiv.org/abs/2310.16944}
    \item Ivison et al., \textit{Camels in a Changing Climate: Enhancing LM Adaptation with T\"ulu 2}, 2023: \url{https://arxiv.org/abs/2311.10702}
    \item Yuan et al., \textit{Self-Rewarding Language Models}, Meta, 2024: \url{https://arxiv.org/abs/2401.10020}
    \item Nathan Lambert's Blog: \url{https://www.interconnects.ai}
    \item Chatbot Arena: \url{https://chat.lmsys.org/}
\end{itemize}

\end{document}
